\documentclass[11pt]{article}
\usepackage[margin=0.9in]{geometry}
\usepackage{fontspec}
\usepackage{amsmath,amssymb,amsthm}
\usepackage{xurl}
\begin{document}\sloppy
\begin{center}
{\Large GROK H0515: cube 30 of \texttt{a40ff54e580f5d0a} closes at 210 of 210 UNSAT; DRAT207; 55 of 75 kept cubes closed}\\[4pt]
2026-10-06 05:15 America/Los\_Angeles \quad branch \texttt{main}, tip \texttt{a980eb6b} \quad automation 41774ea2\\[2pt]
\textbf{P versus NP: OPEN.}
\end{center}

\section*{Peer delta since H0501 (\texttt{9c873cac})}
Checked 05:16\,PT. One commit landed: \texttt{a980eb6b}, ``PROOF-0011: virtual child small-space accounting'', adding \path{research/2026-10-06_0501_PROOF-0011_virtual_child_space.md}. It records VIRTUAL-CHILD-SPACE1 as PROVED/EXACT accounting (simulating an algorithm on a virtual child input by recomputing each requested bit gives $S_{\mathrm{virt}}\le S(M)+N^{o(1)}$ and $T_{\mathrm{virt}}\le T(M)N^{o(1)}$), keeps CSE local addressability and arbitrary-algorithm CSE OPEN, and applies it CONDITIONALLY toward the Williams small-space SAT endpoint with exponent below $2\cos(\pi/7)$; I recomputed $2\cos(\pi/7)=1.80194$ and did not re-read Williams. I read the proof sketch: it is recompute-on-read, the same forget-then-recompute trade that the standing candidate prices with pebbling, and it charges space additively and time multiplicatively, which is consistent with our model. I do not upgrade or downgrade any of its labels. Nothing new under \path{research/main/} or \path{research/agents/grok/}. Root \texttt{main.tex} is not touched.

\section*{1. Literature}
None beyond the peer note above. The next claim is about our own finite search.

\section*{2. Claims recorded}
\textbf{Cube 30, $(0,1,3,2)$, closed (COMPUTATIONAL, solver-trusted).} The splitter log now holds 210 distinct UNSAT subcube indices for cube 30, all 210 with the same step-2 tuple as the matching cube-0 subcube, no SAT line and no \texttt{TIMEOUT} line (the splitter printed \texttt{CUBE-DONE 30} at 05:08:05, which the closure check ignores). Slowest is subcube 134 at 248.2\,s under the 600\,s cap; mean 24.3\,s. The independent closure check counts it closed.

\textbf{Cube 32, $(0,1,3,6)$, partial (COMPUTATIONAL, solver-trusted).} 103 distinct UNSAT subcube indices, 103 of 103 tuples match cube 0, no SAT and no timeout; 107 missing, five running at 05:16\,PT. Slowest so far subcube 74 at 290.6\,s, mean 22.2\,s. Not closed.

\textbf{Claim DRAT207 (COMPUTATIONAL).} 207 of cube 0's 210 subcubes carry drat-trim-verified refutations (new: subcube 52, solve 305.9\,s, check 824.3\,s, \texttt{s VERIFIED}). Remaining 3: subcubes 54, 157, 159; drat-trim on 54 has run since 05:08\,PT under a 3000\,s cap.

\textbf{Status of first-step cubes (COMPUTATIONAL, solver-trusted).} Closed in the closure list: 0, 1, 2, 4, 9, 14, 19, 21, 22, 24, 27, 29, 30. That makes 55 of 75 kept first-step cubes closed, 20 open, 0 SAT (up from 54 and 21 at H0501).

\section*{3. Smallest honest test}
\begin{itemize}
\item \textbf{Closure check} (\path{/workspace/scratch_h0515_1006/closure_run.py}, a copy of H0348's \path{closure.py}; output \path{closure_0516.out}). It ignores \texttt{CUBE-DONE} and needs 210 distinct UNSAT indices with matching tuples. No SAT line; the only timeout is cube 22 subcube 204, superseded by its 3000\,s rerun UNSAT (674.5\,s); closed set as above (13 entries); partial \{32: 100 UNSAT at check time, no open timeout\}.
\item \textbf{Cube scan} (\path{cscan.py 30 32}, output \path{cscan_0516.out}, run seconds later). Cube 30: 210 distinct, 210 of 210 tuples match, 0 missing. Cube 32: 103 distinct, 103 match, 107 missing.
\item \textbf{Recount} (\path{/workspace/scratch_h2333_1005/experiments/recount.py}, rerun 05:16\,PT). Total verified 207, remaining 3, 126 distinct CNF hashes among 126 batch lines, none not verified.
\end{itemize}

\section*{4. What this changes}
Closed rises from 54 to 55 of 75; certified rises from 206 to 207 of 210. Neither moves a charged resource. The bounds stay $9\le L(\texttt{a40ff54e580f5d0a},2)\le 11$ (COMPUTATIONAL for the upper bound; the lower bound stays CONDITIONAL on the tightness lemma as recorded earlier); $L=9$ itself stays open while 20 cubes are open. The largest confirmed price stays one.

\section*{Effect on the standing candidate}
No claim is falsified and no charged resource improves. The time-state-representation surface stays CANDIDATE. PROOF-0011 is accounting that relativizes, by its own audit, so it does not touch the open arrows. The SAT-to-charged-model compiler is OPEN, and so is the product-span obstruction. The finite results concern one fixed function and imply no lower bound for any family.

\section*{Next smallest step}
Let cube 32 finish its last 107 subcubes and run the closure check; rerun any timeout at 3000\,s. Let the DRAT batch finish 54, 157, 159.

\bigskip
\noindent\textbf{Status sentence.} P versus NP remains OPEN: the SAT-to-charged-model compiler and the product-span obstruction are open load-bearing arrows. This pass records that cube 30 of \texttt{a40ff54e580f5d0a} at tight $L=9$ is closed with 210 of 210 subcubes UNSAT and no timeout (COMPUTATIONAL, solver-trusted), that cube 32 is at 103 of 210, and that 207 of cube 0's 210 subcubes carry drat-trim-verified refutations (COMPUTATIONAL); 55 of 75 kept first-step cubes are closed, 20 open, 0 SAT, and $9\le L(f,2)\le 11$ is unchanged. No claim is falsified, no charged resource improves, and this is not a proof of anything about P versus NP.
\end{document}
